\chapter{2002年新课标卷（理）}
\section{单选题}
\begin{problem}
曲线 \(\begin{cases}
x = \cos \theta,\\
y = \sin \theta,
\end{cases}\) （\(\theta\) 为参数）上的点到两坐标轴的距离之和的最大值是（\quad）

\choices
    {\(\frac{1}{2}\)}
    {\(\frac{\sqrt{2}}{2}\)}
    {\(1\)}
    {\(\sqrt{2}\)}
\end{problem}

\begin{answer}
D
\end{answer}

\begin{solution}
曲线上点的坐标满足 \(x=\cos\theta\)，\(y=\sin\theta\)，所以该点到两坐标轴的距离之和为 \(|\cos\theta|+|\sin\theta|\). 由
\((|\cos\theta|+|\sin\theta|)^2=1+2|\sin\theta\cos\theta|\leqslant2\)，得距离之和不超过 \(\sqrt{2}\). 当 \(|\sin\theta|=|\cos\theta|=\frac{\sqrt{2}}{2}\) 时取等号，因此最大值为 \(\sqrt{2}\)，选 D.
\end{solution}

\begin{problem}
复数 \(\left(\frac{1}{2} + \frac{\sqrt{3}}{2}\i\right)^3\) 的值是（\quad）

\choices
    {\(-\i\)}
    {\(\i\)}
    {\(-1\)}
    {\(1\)}
\end{problem}

\begin{answer}
C
\end{answer}

\begin{solution}
因为 \(\frac{1}{2}+\frac{\sqrt{3}}{2}\i=\cos\frac{\pi}{3}+\i\sin\frac{\pi}{3}\)，所以由棣莫弗定理，
\(\left(\frac{1}{2}+\frac{\sqrt{3}}{2}\i\right)^3=\cos\pi+\i\sin\pi=-1\). 因此选 C.
\end{solution}

\begin{problem}
已知 \(m\)，\(n\) 为异面直线，\(m \subset\) 平面 \(\alpha\)，\(n \subset\) 平面 \(\beta\)，\(\alpha \cap \beta = l\)，则 \(l\)（\quad）

\choices
    {与 \(m\)，\(n\) 都相交}
    {与 \(m\)，\(n\) 中至少一条相交}
    {与 \(m\)，\(n\) 都不相交}
    {至多与 \(m\)，\(n\) 中的一条相交}
\end{problem}

\begin{answer}
B
\end{answer}

\begin{solution}
若交线 \(l\) 不与直线 \(m\) 相交，由于 \(l\) 与 \(m\) 同在平面 \(\alpha\) 内，只能有 \(l\parallel m\)；同理，若 \(l\) 不与直线 \(n\) 相交，则 \(l\parallel n\). 若 \(l\) 同时不与 \(m\)，\(n\) 相交，就有 \(m\parallel l\parallel n\)，这与 \(m\)，\(n\) 为异面直线矛盾. 所以 \(l\) 至少与 \(m\)，\(n\) 中的一条相交，选 B.
\end{solution}

\begin{problem}
不等式 \((1 + x)(1 - |x|) > 0\) 的解集是（\quad）

\choices
    {\(\{x\mid 0\leqslant x < 1\}\)}
    {\(\{x\mid x < 0,\ x\neq -1\}\)}
    {\(\{x\mid -1 < x < 1\}\)}
    {\(\{x\mid x < 1,\ x\neq -1\}\)}
\end{problem}

\begin{answer}
D
\end{answer}

\begin{solution}
当 \(x<0\) 时，\(|x|=-x\)，原不等式化为 \((1+x)^2>0\)，所以解为 \(x<0\) 且 \(x\neq-1\). 当 \(x\geqslant0\) 时，原不等式化为 \((1+x)(1-x)>0\)；因 \(1+x>0\)，所以 \(0\leqslant x<1\). 合并两种情况，解集为 \(\{x\mid x<1,\ x\neq-1\}\)，选 D.
\end{solution}

\begin{problem}
在 \((0,2\pi)\) 内，使 \(\sin x > \cos x\) 成立的 \(x\) 的取值范围是（\quad）

\choices
    {\(\left(\frac{\pi}{4},\frac{\pi}{2}\right)\cup \left(\pi,\frac{5\pi}{4}\right)\)}
    {\(\left(\frac{\pi}{4},\pi\right)\)}
    {\(\left(\frac{\pi}{4},\frac{5\pi}{4}\right)\)}
    {\(\left(\frac{\pi}{4},\pi\right) \cup \left(\frac{5\pi}{4},\frac{3\pi}{2}\right)\)}
\end{problem}

\begin{answer}
C
\end{answer}

\begin{solution}
由 \(\sin x>\cos x\) 得 \(\sin x-\cos x>0\)，即 \(\sqrt{2}\sin\left(x-\frac{\pi}{4}\right)>0\). 在 \(0<x<2\pi\) 内，正弦值为正的部分满足 \(0<x-\frac{\pi}{4}<\pi\)，从而 \(\frac{\pi}{4}<x<\frac{5\pi}{4}\). 因此选 C.
\end{solution}

\begin{problem}
设集合 \(M = \left\{x \mid x = \frac{k}{2} + \frac{1}{4}, k \in \Z\right\}\)，\(N = \left\{x \mid x = \frac{k}{4} + \frac{1}{2}, k \in \Z\right\}\)，则（\quad）

\choices
    {\(M = N\)}
    {\(M \subseteq N\)}
    {\(M \supseteq N\)}
    {\(M \cap N = \varnothing\)}
\end{problem}

\begin{answer}
B
\end{answer}

\begin{solution}
集合 \(M\) 中的任意元素可写为 \(x=\frac{2k+1}{4}\). 令 \(j=2k-1\)，则 \(j\in\Z\) 且 \(x=\frac{j}{4}+\frac{1}{2}\)，所以 \(M\subseteq N\). 又 \(\frac{1}{2}\in N\)，而 \(\frac{1}{2}\notin M\)，故 \(M\subset N\) 且 \(M\ne N\)，选 B.
\end{solution}

\begin{problem}
正六棱柱 \( ABCDEF - A_{1}B_{1}C_{1}D_{1}E_{1}F_{1} \) 的底面边长为 \(1\)，侧棱长为 \(\sqrt{2}\)，则这个棱柱侧面对角线 \(E_{1}D\) 与 \(BC_{1}\) 所成的角是（\quad）

\choices
    {\(90^{\circ}\)}
    {\(60^{\circ}\)}
    {\(45^{\circ}\)}
    {\(30^{\circ}\)}
\end{problem}

\begin{answer}
B
\end{answer}

\begin{solution}
在底面内取坐标 \(A(1,0,0)\)，\(B\left(\frac{1}{2},\frac{\sqrt{3}}{2},0\right)\)，\(C\left(-\frac{1}{2},\frac{\sqrt{3}}{2},0\right)\)，\(D(-1,0,0)\)，\(E\left(-\frac{1}{2},-\frac{\sqrt{3}}{2},0\right)\)，并令上底面各点的 \(z\) 坐标为 \(\sqrt{2}\). 则
\(\overrightarrow{E_{1}D}=\left(-\frac{1}{2},\frac{\sqrt{3}}{2},-\sqrt{2}\right)\)，\(\overrightarrow{BC_{1}}=(-1,0,\sqrt{2})\). 因而 \(\overrightarrow{E_{1}D}\cdot\overrightarrow{BC_{1}}=-\frac{3}{2}\)，且 \(|E_{1}D|=|BC_{1}|=\sqrt{3}\). 设两直线所成的锐角为 \(\theta\)，则
\(\cos\theta=\frac{|\overrightarrow{E_{1}D}\cdot\overrightarrow{BC_{1}}|}{|E_{1}D|\,|BC_{1}|}=\frac{1}{2}\)，所以 \(\theta=60^{\circ}\)，选 B.
\end{solution}

\begin{problem}
函数 \(y = x^{2} + bx + c\)（\(x \in [0, +\infty)\)）是单调函数的充要条件是（\quad）

\choices
    {\(b \geqslant 0\)}
    {\(b \leqslant 0\)}
    {\(b > 0\)}
    {\(b < 0\)}
\end{problem}

\begin{answer}
A
\end{answer}

\begin{solution}
任取 \(0\leqslant x<y\)，有 \(f(y)-f(x)=(y-x)(x+y+b)\). 若 \(b\geqslant0\)，则 \(x+y+b>0\)，所以 \(f(y)>f(x)\)，函数在 \([0,+\infty)\) 上单调递增. 若 \(b<0\)，取 \(x=0\)，\(y=-\frac{b}{2}\)，则 \(f(y)-f(0)=y(y+b)=-\frac{b^{2}}{4}<0\)；再取 \(x=-b\)，\(y=-b+1\)，则 \(f(y)-f(x)=1-b>0\). 因此 \(b<0\) 时函数既不单调递增也不单调递减. 所以充要条件是 \(b\geqslant0\)，选 A.
\end{solution}

\begin{problem}
已知 \(0 < x < y < a < 1\)，则有（\quad）

\choices
    {\(\log_a(xy) < 0\)}
    {\(0 < \log_a(xy) < 1\)}
    {\(1 < \log_a(xy) < 2\)}
    {\(\log_a(xy) > 2\)}
\end{problem}

\begin{answer}
D
\end{answer}

\begin{solution}
由 \(0<x<a\)，\(0<y<a\)，得 \(0<xy<a^{2}\). 因为 \(0<a<1\)，函数 \(\log_{a}t\) 在正数范围内递减，所以
\(\log_{a}(xy)>\log_{a}(a^{2})=2\). 因此选 D.
\end{solution}

\begin{problem}
平面直角坐标系中，\(O\) 为坐标原点，已知两点 \(A(3,1)\)，\(B(-1,3)\)，若点 \(C\) 满足 \(\overrightarrow{OC} = \alpha \overrightarrow{OA} + \beta \overrightarrow{OB}\)，其中 \(\alpha\)，\(\beta \in \R\)，且 \(\alpha + \beta = 1\)，则点 \(C\) 的轨迹方程为（\quad）

\choices
    {\(3x + 2y - 11 = 0\)}
    {\((x - 1)^{2} + (y - 2)^{2} = 5\)}
    {\(2x - y = 0\)}
    {\(x + 2y - 5 = 0\)}
\end{problem}

\begin{answer}
D
\end{answer}

\begin{solution}
由 \(\alpha+\beta=1\) 得 \(\beta=1-\alpha\). 设 \(C(x,y)\)，则
\(x=3\alpha-\beta=4\alpha-1\)，\(y=\alpha+3\beta=3-2\alpha\). 消去 \(\alpha\)，得 \(x+2y-5=0\). 由于 \(\alpha\) 可取任意实数，轨迹是该直线，选 D.
\end{solution}

\begin{problem}
从正方体的 \(6\text{ 个}\)面中选取 \(3\text{ 个}\)面，其中有 \(2\text{ 个}\)面不相邻的选法共有（\quad）

\choices
    {\(8\text{ 种}\)}
    {\(12\text{ 种}\)}
    {\(16\text{ 种}\)}
    {\(20\text{ 种}\)}
\end{problem}

\begin{answer}
B
\end{answer}

\begin{solution}
正方体的面可分成 \(3\) 对相对的面. 若所选的 \(3\) 个面中有一对不相邻的面，它们必须是一对相对面；先选这对面有 \(3\) 种选法，再从剩下的 \(4\) 个面中选第 \(3\) 个面. 三个面中不可能同时出现两对相对面，所以选法总数为 \(3\times4=12\). 因此选 B.
\end{solution}

\begin{problem}
据 \(2002\text{ 年}\) \(3\text{ 月}\) \(5\text{ 日}\)九届人大五次会议《政府工作报告》：“\(2001\text{ 年}\)国内生产总值达到 \(95933\text{ 亿元}\)，比上年增长 \(7.3\%\)”，如果“十·五”期间（\(2001\text{ 年}\)－\(2005\text{ 年}\)）每年的国内生产总值都按此年增长率增长，那么到“十·五”末我国国内年生产总值约为（\quad）

\choices
    {\(115000\text{ 亿元}\)}
    {\(120000\text{ 亿元}\)}
    {\(127000\text{ 亿元}\)}
    {\(135000\text{ 亿元}\)}
\end{problem}

\begin{answer}
C
\end{answer}

\begin{solution}
从 \(2001\) 年到 \(2005\) 年共有 \(4\) 次按增长率增长，所以末年的国内生产总值为
\(95933(1+7.3\%)^{4}=95933\times1.073^{4}\). 计算得 \(1.073^{2}=1.151329\)，从而 \(1.073^{4}\approx1.3256\)，故总值约为 \(95933\times1.3256\approx127000\) 亿元. 因此选 C.
\end{solution}

\section{填空题}
\begin{problem}
函数 \(y = \frac{2x}{1 + x}\)（\(x \in (-1, +\infty)\)）图象与其反函数图象的交点为\fillinblank{}.
\end{problem}

\begin{answer}
\((0,0)\)，\((1,1)\)
\end{answer}

\begin{solution}
由 \(y=\frac{2x}{1+x}\) 解得反函数为 \(y=\frac{x}{2-x}\)，其定义域为 \(x<2\). 设交点为 \((u,v)\)，则
\(v=\frac{2u}{1+u}\)，且 \(v=\frac{u}{2-u}\). 交点的横坐标必须同时满足原函数与反函数的定义域，所以 \(u>-1\) 且 \(u<2\). 两式相等得
\(\frac{2u}{1+u}=\frac{u}{2-u}\)，从而 \(2u(2-u)=u(1+u)\)，即 \(u(3-3u)=0\). 所以 \(u=0\) 或 \(u=1\). 代回 \(v=\frac{2u}{1+u}\)，分别得到 \(v=0\) 或 \(v=1\)，且两点均满足定义域条件，故交点为 \((0,0)\)，\((1,1)\).
\end{solution}

\begin{problem}
椭圆 \(5x^{2} + ky^{2} = 5\) 的一个焦点是 \((0, 2)\)，那么 \(k =\)\fillinblank{}.
\end{problem}

\begin{answer}
1
\end{answer}

\begin{solution}
将椭圆方程除以 \(5\)，得 \(x^{2}+\frac{k}{5}y^{2}=1\). 因为焦点 \((0,2)\) 在 \(y\) 轴上，设纵向半轴为 \(b\)，则横向半轴为 \(1\)，且焦半距 \(c=2\). 由椭圆关系 \(c^{2}=b^{2}-1^{2}\)，得 \(b^{2}=5\). 又 \(b^{2}=\frac{5}{k}\)，所以 \(k=1\).
\end{solution}

\begin{problem}
直线 \(x = 0\)，\(y = 0\)，\(x = 2\) 与曲线 \(y = (\sqrt{2})^{x}\) 所围成的图形绕 \(x\) 轴旋转一周而成的旋转体的体积等于\fillinblank{}.
\end{problem}

\begin{answer}
\(\frac{3\pi}{\ln 2}\)
\end{answer}

\begin{solution}
曲线为 \(y=(\sqrt{2})^{x}=2^{x/2}\)，在 \(0\leqslant x\leqslant2\) 内，所围图形绕 \(x\) 轴旋转所得体积为
\(V=\pi\int_{0}^{2}y^{2}\,\mathrm{d}x=\pi\int_{0}^{2}2^{x}\,\mathrm{d}x=\frac{\pi}{\ln2}\left.2^{x}\right|_{0}^{2}=\frac{3\pi}{\ln2}\).
\end{solution}

\begin{problem}
已知 \(f(x) = \frac{x^2}{1 + x^2}\)，那么 \(f(1) + f(2) + f\left(\frac{1}{2}\right) + f(3) + f\left(\frac{1}{3}\right) + f(4) + f\left(\frac{1}{4}\right) =\)\fillinblank{}.
\end{problem}

\begin{answer}
\(\frac{7}{2}\)
\end{answer}

\begin{solution}
对任意非零的 \(t\)，有 \(f(t)+f\left(\frac{1}{t}\right)=\frac{t^{2}}{1+t^{2}}+\frac{1}{1+t^{2}}=1\). 因而
\(f(2)+f\left(\frac{1}{2}\right)=1\)，\(f(3)+f\left(\frac{1}{3}\right)=1\)，\(f(4)+f\left(\frac{1}{4}\right)=1\)，且 \(f(1)=\frac{1}{2}\). 所求和为 \(1+1+1+\frac{1}{2}=\frac{7}{2}\).
\end{solution}

\section{解答题}
\begin{problem}
已知 \(\cos \left( \alpha + \frac{\pi}{4} \right) = \frac{3}{5}\)，\(\frac{\pi}{2} \leqslant \alpha < \frac{3\pi}{2}\)，求 \(\cos \left(2\alpha + \frac{\pi}{4}\right)\) 的值.
\end{problem}

\begin{solution}
令 \(\beta=\alpha+\frac{\pi}{4}\)，则 \(\cos\beta=\frac{3}{5}\)，且 \(\frac{3\pi}{4}\leqslant\beta<\frac{7\pi}{4}\). 在此区间内余弦为正只能发生在第四象限，所以 \(\sin\beta=-\frac{4}{5}\). 因此
\(\cos2\beta=\cos^{2}\beta-\sin^{2}\beta=\frac{9}{25}-\frac{16}{25}=-\frac{7}{25}\)，\(\sin2\beta=2\sin\beta\cos\beta=-\frac{24}{25}\). 又 \(2\alpha+\frac{\pi}{4}=2\beta-\frac{\pi}{4}\)，故
\(\cos\left(2\alpha+\frac{\pi}{4}\right)=\cos\left(2\beta-\frac{\pi}{4}\right)=\frac{\sqrt{2}}{2}\left(\cos2\beta+\sin2\beta\right)=-\frac{31\sqrt{2}}{50}\).
\end{solution}

\section{选做题：解答题}
\begin{problem}[甲]
如图，正三棱柱 \(ABC - A_{1}B_{1}C_{1}\) 的底面边长为 \(a\)，侧棱长为 \(\sqrt{2}a\).

\begin{enumerate}
\item 建立适当的坐标系，并写出点 \(A\)，\(B\)，\(A_{1}\)，\(C_{1}\) 的坐标；
\item 求 \(AC_{1}\) 与侧面 \(ABB_{1}A_{1}\) 所成的角.
\end{enumerate}

\bitmapfigure[max width=0.36\linewidth]{img/2002/new_curriculum_science/q18_fig1.png}
\FloatBarrier
\end{problem}

\begin{solution}
\begin{enumerate}
\item 在底面内取 \(xOy\) 平面，使 \(A(0,0,0)\)，\(B(0,a,0)\)，\(C\left(-\frac{\sqrt{3}}{2}a,\frac{a}{2},0\right)\)，再取 \(z\) 轴垂直于底面. 则 \(A_{1}(0,0,\sqrt{2}a)\)，\(C_{1}\left(-\frac{\sqrt{3}}{2}a,\frac{a}{2},\sqrt{2}a\right)\).
\item 侧面 \(ABB_{1}A_{1}\) 的方程为 \(x=0\)，其一个法向量为 \((1,0,0)\). 直线 \(AC_{1}\) 的方向向量为 \(\overrightarrow{AC_{1}}=\left(-\frac{\sqrt{3}}{2}a,\frac{a}{2},\sqrt{2}a\right)\)，且 \(|AC_{1}|=\sqrt{3}a\). 设直线与平面所成的角为 \(\theta\)，则
\(\sin\theta=\frac{|\overrightarrow{AC_{1}}\cdot(1,0,0)|}{|AC_{1}|}=\frac{\frac{\sqrt{3}}{2}a}{\sqrt{3}a}=\frac{1}{2}\)，所以 \(\theta=30^{\circ}\).
\end{enumerate}
\end{solution}

\reuseproblemnumber
\begin{problem}[乙]
如图，正方形 \(ABCD\)，\(ABEF\) 的边长都是 \(1\)，而且平面 \(ABCD\)，\(ABEF\) 互相垂直. 点 \(M\) 在 \(AC\) 上移动，点 \(N\) 在 \(BF\) 上移动，若 \(CM = BN = a\)（\(0 < a < \sqrt{2}\)）.

\begin{enumerate}
\item 求 \(MN\) 的长；
\item \(a\) 为何值时，\(MN\) 的长最小；
\item 当 \(MN\) 的长最小时，求面 \(MNA\) 与面 \(MNB\) 所成二面角 \(\alpha\) 的大小.
\end{enumerate}

\bitmapfigure[max width=0.36\linewidth]{img/2002/new_curriculum_science/q19_fig1.png}
\FloatBarrier
\end{problem}

\begin{solution}
\begin{enumerate}
\item 取空间直角坐标系，使 \(A(0,0,0)\)，\(B(1,0,0)\)，\(C(1,1,0)\)，\(D(0,1,0)\)，\(E(1,0,1)\)，\(F(0,0,1)\). 于是正方形 \(ABCD\) 在 \(z=0\) 内，正方形 \(ABEF\) 在 \(y=0\) 内，且两平面互相垂直. 因为 \(AC=BF=\sqrt{2}\)，由 \(CM=BN=a\) 得
\(M\left(1-\frac{\sqrt{2}}{2}a,1-\frac{\sqrt{2}}{2}a,0\right)\)，\(N\left(1-\frac{\sqrt{2}}{2}a,0,\frac{\sqrt{2}}{2}a\right)\). 因而
\(MN^{2}=\left(1-\frac{\sqrt{2}}{2}a\right)^{2}+\left(\frac{\sqrt{2}}{2}a\right)^{2}=\left(a-\frac{\sqrt{2}}{2}\right)^{2}+\frac{1}{2}\)，所以 \(MN=\sqrt{\left(a-\frac{\sqrt{2}}{2}\right)^{2}+\frac{1}{2}}\).
\item 由上式，\(MN\) 取最小值当且仅当 \(a-\frac{\sqrt{2}}{2}=0\). 由于 \(0<a<\sqrt{2}\)，该值符合条件，所以当 \(a=\frac{\sqrt{2}}{2}\) 时，\(MN\) 的最小值为 \(\frac{\sqrt{2}}{2}\).
\item 此时 \(M\left(\frac{1}{2},\frac{1}{2},0\right)\)，\(N\left(\frac{1}{2},0,\frac{1}{2}\right)\). 取
\(\bs{n}_{1}=\overrightarrow{MA}\times\overrightarrow{MN}=\frac{1}{4}(-1,1,1)\)，\(\bs{n}_{2}=\overrightarrow{MB}\times\overrightarrow{MN}=\frac{1}{4}(-1,-1,-1)\)，它们分别是平面 \(MNA\)，\(MNB\) 的法向量. 于是
\(\cos\alpha=\frac{\bs{n}_{1}\cdot\bs{n}_{2}}{|\bs{n}_{1}|\,|\bs{n}_{2}|}=-\frac{1}{3}\)，故 \(\alpha=\arccos\left(-\frac{1}{3}\right)\).
\end{enumerate}
\end{solution}

\section{解答题}
\begin{problem}
某单位 \(6\text{ 个}\)员工借助互联网开展工作，每个员工上网的概率都是 \(0.5\)（相互独立）.

\begin{enumerate}
\item 求至少 \(3\text{ 人}\)同时上网的概率；
\item 至少几人同时上网的概率小于 \(0.3\)？
\end{enumerate}
\end{problem}

\begin{solution}
\begin{enumerate}
\item 设同时上网的人数为 \(X\)，则 \(X\sim B\left(6,\frac{1}{2}\right)\). 至少 \(3\) 人同时上网包括恰有 \(3\)，\(4\)，\(5\)，\(6\) 人同时上网，所以
\(P(X\geqslant3)=\sum_{k=3}^{6}\mathrm C_{6}^{k}\left(\frac{1}{2}\right)^{k}\left(\frac{1}{2}\right)^{6-k}=\frac{\mathrm C_{6}^{3}+\mathrm C_{6}^{4}+\mathrm C_{6}^{5}+\mathrm C_{6}^{6}}{2^{6}}=\frac{20+15+6+1}{64}=\frac{21}{32}\).
\item 因为 \(P(X\geqslant4)=\frac{15+6+1}{64}=\frac{22}{64}=0.34375>0.3\)，而 \(P(X\geqslant5)=\frac{6+1}{64}=\frac{7}{64}=0.109375<0.3\)，所以至少 \(5\) 人同时上网时，该概率小于 \(0.3\)，且 \(5\) 是满足条件的最小人数.
\end{enumerate}
\end{solution}

\begin{problem}
已知 \(a > 0\)，函数 \(f(x) = \frac{1 - ax}{x}\)（\(x \in (0, +\infty)\)）. 设 \(0 < x_1 < \frac{2}{a}\). 记曲线 \(y = f(x)\) 在点 \(M(x_1, f(x_1))\) 处的切线为 \(l\).

\begin{enumerate}
\item 求 \(l\) 的方程；
\item 设 \(l\) 与 \(x\) 轴交点为 \((x_{2},0)\). 证明：
\begin{enumerate}
\item \(0 < x_{2} \leqslant \frac{1}{a}\)；
\item 若 \(x_{1} < \frac{1}{a}\)，则 \(x_{1} < x_{2} < \frac{1}{a}\).
\end{enumerate}
\end{enumerate}
\end{problem}

\begin{solution}
\begin{enumerate}
\item \(f(x)=\frac{1-ax}{x}=\frac{1}{x}-a\)，所以 \(f'(x)=-\frac{1}{x^{2}}\). 点 \(M\) 的坐标为 \(\left(x_{1},\frac{1}{x_{1}}-a\right)\)，故切线 \(l\) 的方程为
\(y-\left(\frac{1}{x_{1}}-a\right)=-\frac{1}{x_{1}^{2}}(x-x_{1})\).

\item 令切线与 \(x\) 轴交于 \((x_{2},0)\)，代入上式得
\(0-\left(\frac{1}{x_{1}}-a\right)=-\frac{1}{x_{1}^{2}}(x_{2}-x_{1})\)，从而 \(x_{2}=x_{1}(2-ax_{1})\).
\begin{enumerate}
\item 由 \(0<x_{1}<\frac{2}{a}\) 得 \(0<ax_{1}<2\)，所以 \(x_{2}=x_{1}(2-ax_{1})>0\). 又
\(ax_{2}=ax_{1}(2-ax_{1})=1-(ax_{1}-1)^{2}\leqslant1\)，故 \(0<x_{2}\leqslant\frac{1}{a}\).
\item 若 \(x_{1}<\frac{1}{a}\)，则 \(0<ax_{1}<1\). 因而
\(x_{2}-x_{1}=x_{1}(1-ax_{1})>0\)，即 \(x_{2}>x_{1}\). 同时 \(1-ax_{2}=(1-ax_{1})^{2}>0\)，即 \(x_{2}<\frac{1}{a}\). 所以 \(x_{1}<x_{2}<\frac{1}{a}\).
\end{enumerate}
\end{enumerate}
\end{solution}

\begin{problem}
已知两点 \(M(-1,0)\)，\(N(1,0)\). 且点 \(P\) 使 \(\overrightarrow{MP} \cdot \overrightarrow{MN}\)，\(\overrightarrow{PM} \cdot \overrightarrow{PN}\)，\(\overrightarrow{NM} \cdot \overrightarrow{NP}\) 成公差小于零的等差数列.

\begin{enumerate}
\item 点 \(P\) 的轨迹是什么曲线？
\item 若点 \(P\) 坐标为 \((x_0, y_0)\)，记 \(\theta\) 为 \(\overrightarrow{PM}\) 与 \(\overrightarrow{PN}\) 的夹角，求 \(\tan \theta\).
\end{enumerate}
\end{problem}

\begin{solution}
\begin{enumerate}
\item 设 \(P(x,y)\). 则
\(\overrightarrow{MP}=(x+1,y)\)，\(\overrightarrow{MN}=(2,0)\)，\(\overrightarrow{PM}=(-x-1,-y)\)，\(\overrightarrow{PN}=(1-x,-y)\)，\(\overrightarrow{NM}=(-2,0)\)，\(\overrightarrow{NP}=(x-1,y)\). 因此三个数量积依次为
\(\overrightarrow{MP}\cdot\overrightarrow{MN}=2(x+1)\)，\(\overrightarrow{PM}\cdot\overrightarrow{PN}=x^{2}+y^{2}-1\)，\(\overrightarrow{NM}\cdot\overrightarrow{NP}=2(1-x)\).

三数成等差数列的充要条件是首项与末项之和等于中项的两倍，所以
\(2(x+1)+2(1-x)=2(x^{2}+y^{2}-1)\)，即 \(x^{2}+y^{2}=3\). 其公差为
\((x^{2}+y^{2}-1)-2(x+1)=-2x\). 由公差小于零得 \(x>0\). 反之，圆上任一点满足 \(x>0\) 时，上述三数公差均为 \(-2x<0\)，所以点 \(P\) 的轨迹为 \(x^{2}+y^{2}=3\)，\(x>0\) 的部分，即不含端点的右半圆.
\item 在轨迹上 \(x_{0}^{2}+y_{0}^{2}=3\)，所以
\(\overrightarrow{PM}\cdot\overrightarrow{PN}=x_{0}^{2}+y_{0}^{2}-1=2\). 由两向量的二维叉积的绝对值，
\(\left|\overrightarrow{PM}\times\overrightarrow{PN}\right|=\left|(-x_{0}-1)(-y_{0})-(-y_{0})(1-x_{0})\right|=2|y_{0}|\). 因为数量积为正，\(\theta\) 为锐角或零角，故
\(\tan\theta=\frac{2|y_{0}|}{2}=|y_{0}|=\sqrt{3-x_{0}^{2}}\).
\end{enumerate}
\end{solution}

\begin{problem}
已知 \(\{a_{n}\}\) 是由非负整数组成的数列，满足 \(a_{1} = 0\)，\(a_{2} = 3\)，\(a_{n+1}a_{n} = (a_{n-1} + 2)(a_{n-2} + 2)\)，\(n = 3\)，\(4\)，\(5\)，\(\cdots\).

\begin{enumerate}
\item 求 \(a_{3}\)；
\item 证明 \(a_{n} = a_{n - 2} + 2\)，\(n = 3\)，\(4\)，\(5\)，\(\cdots\)；
\item 求 \(\{a_{n}\}\) 的通项公式及其前 \(n\) 项和 \(S_{n}\).
\end{enumerate}
\end{problem}

\begin{solution}
\begin{enumerate}
\item 令 \(a_{3}=p\). 在递推式中取 \(n=3\)，得 \(a_{4}p=(a_{2}+2)(a_{1}+2)=10\). 因为各项是非负整数，且右端为正，所以 \(p\) 是 \(10\) 的正整数因数. 取 \(n=4\)，有
\(a_{5}a_{4}=(a_{3}+2)(a_{2}+2)=5(p+2)\)，从而 \(a_{5}=\frac{p(p+2)}{2}\). 当 \(p=1\) 或 \(p=5\) 时，\(a_{5}\) 分别为 \(\frac{3}{2}\) 或 \(\frac{35}{2}\)，不合整数条件；当 \(p=2\) 时，\(a_{4}=5\)，\(a_{5}=4\). 若 \(p=10\)，则 \(a_{4}=1\)，\(a_{5}=60\)，再取 \(n=5\) 得 \(60a_{6}=(1+2)(10+2)=36\)，这使 \(a_{6}=\frac{3}{5}\)，也不合整数条件. 所以 \(a_{3}=2\).
\item 由 \(a_{3}=2=a_{1}+2\)，且 \(a_{4}=5=a_{2}+2\)，关系 \(a_{n}=a_{n-2}+2\) 在 \(n=3\)，\(4\) 时成立. 假设在某个 \(n\geqslant3\) 时 \(a_{n}=a_{n-2}+2\)，则 \(a_{n}>0\)，并由递推式得
\(a_{n+1}a_{n}=(a_{n-1}+2)(a_{n-2}+2)=(a_{n-1}+2)a_{n}\). 因为 \(a_{n}>0\)，可约去 \(a_{n}\)，得到 \(a_{n+1}=a_{n-1}+2\). 因此由数学归纳法，对一切 \(n=3\)，\(4\)，\(5\)，\(\cdots\)，都有 \(a_{n}=a_{n-2}+2\).
\item 奇数项和偶数项分别为等差数列：\(a_{2k-1}=2(k-1)=2k-2\)，\(a_{2k}=3+2(k-1)=2k+1\). 因而统一写成 \(a_{n}=n+(-1)^{n}\). 于是
\(S_{n}=\sum_{k=1}^{n}a_{k}=\sum_{k=1}^{n}k+\sum_{k=1}^{n}(-1)^{k}\). 当 \(n\) 为偶数时，第二个和为 \(0\)；当 \(n\) 为奇数时，第二个和为 \(-1\). 所以
\[
S_{n}=\begin{cases}
\dfrac{1}{2}n(n+1),& n\text{ 为偶数},\\
\dfrac{1}{2}n(n+1)-1,& n\text{ 为奇数}.
\end{cases}
\]
\end{enumerate}
\end{solution}
